\section{Nonlinearity and the Product Block}
\label{sec:nonlinear_product}

% Arc: nonlinearity gets its own section. Tree results open it — the tree edge
% is the tease that motivates the product block of the latest model. Ends with
% the 3-block ridge: alpha=1 target-HAR + alpha=100 linear exog block +
% alpha=1000 product block.
% Base text: drafts_2026-07-30/linear_vs_nonlinear.tex (nonlinear half),
% reference/sections/linear_vs_nonlinear.tex (tree premium, saturation at
% additive-plus-pairwise anchored to the session clock).

Section~\ref{sec:dense_weak} closed the linear question: the exogenous signal
is dense but weak, uniform shrinkage is the correct linear response, and the
two-block ridge is where that response ends. This section asks what a richer
hypothesis class finds that the linear one cannot. The answer comes in three
movements. First, gradient-boosted trees hold a real, uniform, but modest edge
over the best linear arm --- the one model class in the study that beats the
shrunk linear reader. Second, when that edge is decomposed it turns out to be
shallow and nameable: almost all of it is additive-plus-pairwise structure,
and the pairwise part is a small set of \emph{products} of features the panel
already contains, anchored to the session clock. Third --- and this is the
constructive payoff --- once the structure is named it does not need a tree.
The products can be written down as explicit columns, handed to the same
rolling ridge machinery as everything else, and shrunk as a block of their
own. The section ends with the resulting three-block ridge, the second rung of
the block ladder this paper builds (two blocks in
Section~\ref{sec:dense_weak}, four in Section~\ref{sec:transmission}).

\subsection{The Tree Edge}
\label{sec:tree_edge}

\paragraph{The comparison.} The evidence comes from the causal-tuning battery
of Section~\ref{sec:buckets}: five estimators fully crossed with the feature
buckets on the shared panel ($n = 218{,}934$ thirty-minute bars, identical
rows in every cell), every arm re-selecting its own hyperparameters
periodically and causally, so that no forecast's configuration ever sees its
own future. The two nonlinear arms are the gradient-boosted tree ensembles
LightGBM \citep{Ke2017} and XGBoost \citep{Chen2016}; the linear arms are the
penalized regressions of Section~\ref{sec:dense_weak}; the shared comparator
is the per-bar OLS--HAR incumbent at QLIKE $0.13415$
(Section~\ref{sec:dp_benchmark}). Because the tree arms are causally tuned by
the same protocol as the linear arms, the comparison measures function class,
not search effort.

The ordering is uniform. LightGBM is the best arm on every one of the nine
feature sets; the best linear arm (ridge) is second on every one; there is no
bucket on which the ranking reverses. On the richest set,
\texttt{all\_features}, LightGBM reaches $0.12604$ against ridge's $0.12788$
--- a model-class gap of $0.00184$ --- and the tree-minus-linear difference
across the nine sets ranges from $0.00184$ to $0.00358$. An exact sign test
over the nine sets gives $p = 0.0039$, and a bound-based pairwise test
(the incumbent cancels from the tree-minus-linear loss differential, so the
pairwise standard error is bounded by the sum of the two arms' standard errors
without any assumption on their dependence) resolves all nine matched
comparisons in the tree's favour after Holm adjustment: the headline
\texttt{all\_features} pair is the weakest at $|t| \ge 2.19$, $p \le 0.029$,
the other eight ranging from $|t| \ge 3.73$ to $|t| \ge 5.91$. The edge is
small --- roughly a quarter of the information effect of
Section~\ref{sec:buckets} --- but it is real, and it is systematic.

\paragraph{The pure-nonlinearity premium.} The battery's no-exogenous
\texttt{baseline} row isolates function class exactly, because the information
set in that row is identical to the incumbent's: HAR lags and calendar
features, nothing else. There the causally tuned LightGBM reaches $0.13118$
against the incumbent's $0.13415$ ($t = -5.6$; XGBoost $0.13163$, $t = -5.5$)
--- a premium of $0.00297$ available from relaxing the linear functional form
alone. Two corrections temper that headline. Measured against the base-2 OLS
ladder at $0.13332$ rather than the quadrature-limited base-5 incumbent, the
premium shrinks to $0.00215$: roughly $28\%$ of the naive figure is HAR
lag-grid quadrature error that a tree on a log-spaced boxcar basis
interpolates away, not nonlinearity. And the comparison sets a causally tuned
tree against an untuned OLS; against the causally tuned ridge on the same
features ($0.13445$) the premium is $0.00327$. We carry the conservative
$0.00215$ forward. Meanwhile the three linear arms on that same row are the
only three arms in the battery that \emph{lose} to the incumbent --- shrinkage
with nothing to shrink is pure bias --- so on the HAR-plus-calendar block the
function-class question and the penalty question separate cleanly: penalties
do nothing, and the tree finds something real.

\paragraph{The signature of how the tree wins.} The Diebold--Mariano
statistics \citep{DieboldMariano1995,HarveyLeybourneNewbold1997} carry
information the QLIKE levels do not, and on \texttt{all\_features} they invert
the ranking: ridge, the worse model by level, achieves $t = -19.9$ against the
incumbent, while LightGBM, the better model, achieves $t = -15.4$. A larger
$|t|$ at a smaller mean improvement implies a smaller bar-to-bar dispersion of
that improvement: the linear model wins by a little on almost every bar,
while the tree wins by more on average but with substantially more
variability. That is what a localized correction looks like --- the tree
leaves most of the state space to its additive base and intervenes in
particular regions --- and it is the first hint that the tree's edge is a
\emph{specific, findable structure} rather than a diffuse capacity advantage.

\paragraph{Where the nonlinearity lives.} Three independent decompositions,
carried over from the earlier phase of this project, agree on what that
structure is.

\emph{It is shallow.} A functional-ANOVA decomposition of a fitted depth-eight
boosted tree attributes approximately $55\%$ of its explainable variance to
additive terms, $9\%$ to pairwise interactions, and $36\%$ to a diffuse
higher-order remainder concentrated in no identifiable interaction; within the
additive part the HAR block alone accounts for $96\%$. A model nominally
capable of eighth-order interactions spends its capacity on main effects plus
a thin pairwise layer.

\emph{It saturates at pairwise.} On the deployed residual-stack panel of the
earlier build ($n = 194{,}934$; levels not comparable to the battery), an
expressivity ladder that holds the base model fixed and varies only the
correction stage closes the gap monotonically and completes the closing at
pairwise order: a purely additive stage is inadequate at $0.12757$, adding
bilinear pairwise interactions recovers nearly everything at $0.12108$,
nonlinear pairwise shapes reach $0.12080$, and the deployed bagged
additive-plus-pairwise stage (a GA2M in the sense of \citealt{Lou2013}) is
best at $0.12033$ --- while an Optuna-searched unrestricted-depth XGBoost
stage lands at $0.12094$, \emph{worse than fitting no correction at all}
($0.12081$). Unrestricted capacity, properly searched, is a net negative;
restricted pairwise capacity is the whole story.

\emph{It distills into products.} Most decisively, the nonlinearity is
substantially recoverable in closed form, and the closed form is a
\emph{product of existing columns}. On an earlier build of the same panel
family, a plain elastic-net base scores $0.12516$ and a residualized
boosted-tree correction takes it to $0.12414$; adding a single hand-written
interaction --- the HAR lags multiplied by a late-day session dummy, six
columns --- reaches $0.12457$, recovering $58\%$ of the tree's entire edge,
and the full set of HAR-by-six-regime interactions reaches $0.12453$, or
$62\%$. What the tree discovers is not a rich functional manifold; it is that
HAR's persistence coefficient changes sign near the session close (the
clock-anchored sign-flip documented descriptively in
Section~\ref{sec:data_prep}), and it encodes that fact through splits because
splits are the only vocabulary it has. A linear model given the product term
says the same thing directly.

The conclusion of the tree experiments is therefore not ``use trees.'' It is
a specification hint: \textbf{the panel's recoverable nonlinearity is a
sparse set of pairwise products of columns the design already contains}. The
rest of this section takes the hint literally.

\subsection{The Product Block}
\label{sec:product_block}

If the tree's edge is pairwise products, the natural move is to enumerate
candidate products, select the few that matter, and append them to the linear
design as ordinary columns --- turning the tree's implicit interactions into
an explicit feature block that the same rolling ridge machinery can read. The
interaction study that developed this block answered four questions:
which products, chosen how often, shrunk how hard, refit how fast. Its
numbers were computed on the interaction-study panel --- $218{,}934$
out-of-sample bars, 2007--2024, a $250$-day rolling window, HAR reference
QLIKE $0.13275$ --- which shares its rows with the battery panel but not its
window or backbone convention, so levels below are not comparable to the
battery tables and are quoted only for the design decisions they settle.
\pending{confirmation that the paper's production product block is exactly
this construction (candidate space, selection window, column scaling) once
the block is re-run under the paper's panel convention}

\paragraph{The candidate space, and selection.} Candidates are the pairwise
products (including squares) of $133$ base columns --- the HAR ladder, the
$41$ exogenous channels at three of their rolling windows, and the clock
gates --- giving $8{,}911$ candidate products. Selection is by marginal
correlation with the forecast residual, and the selected-set-size ladder is
cleanly unimodal: accuracy peaks at $k = 25$--$100$ of $8{,}911$ and degrades
by $k = 400$. This is the mirror image of the linear channel of
Section~\ref{sec:dense_weak}, where accuracy improved monotonically in the
number of columns: \emph{the linear channel is dense; the interaction channel
is sparse}. Roughly a hundred products carry what the trees were finding.

\paragraph{Frozen, not tracked.} The sharpest design fact is temporal. An arm
that reselects the product set every month --- identical machinery, identical
coefficient refits, only the selection differing --- \emph{loses} to an arm
that selects the $100$ products once, from the 2005--2006 window, and never
touches them again: at $k = 100$ the dynamic arm trails the static one at
DM $t = -3.21$ under the study's original scaling, and the sign is negative
at every healthy specification tested. The dynamic set churns $24$--$31\%$ of
its members per month, and that churn is noise that costs real accuracy to
chase. The product block is therefore \emph{sparse in identity and static in
time}: which interactions matter was decided once, eighteen years before the
end of the sample, and reselection never pays.

\paragraph{The products are nameable.} The persistently selected products are
not opaque. Nine of the top ten involve the signed bar return: its square (a
variance-ratio/jump term), signed return $\times$ volume (order-flow
imbalance in all but name), signed return $\times$ the HAR level (the
leverage effect conditional on the volatility state), signed return $\times$
absolute return (a semivariance-style asymmetry), signed return $\times$
fourth-moment activity. The hub is the signed return, whose predictive value
is almost entirely interactive rather than marginal --- a ridge can form any
linear combination of the panel's columns, but it structurally cannot form a
product, and these products are precisely the asymmetry and flow terms the
volatility literature would nominate. The block is the tree's edge with
names attached.

\paragraph{The block needs its own penalty.} Products of centered features
have variance on the order of the product of their parents' variances, so a
product block and a base block live on incommensurate scales, and one global
ridge penalty cannot restrain the former while fitting the latter. The
interaction study learned this the hard way: under a single shared penalty
the product block's contribution is unstable and can be outright negative,
and an early $+43\%$ headline for the block turned out to be an artifact of a
design clip that was silently standing in for the missing block-wise
regularization. Under an explicit separate penalty the picture is clean and
monotone: with the product block shrunk an order of magnitude harder than the
linear block, the hundred frozen products add $\Delta R^2 = +0.0067$ over the
full $516$-column linear base --- a $+35\%$ relative improvement in
residual-space accuracy, DM $t = +2.71$ --- while lighter product penalties
destroy the gain entirely. The block-penalty ratio is not a tuning nicety; it
is the difference between the interaction channel helping and hurting.

\paragraph{The block needs fresh coefficients.} The products' value also
depends on refit cadence more strongly than the linear block's does. At the
monthly coefficient refit the block inherited from the selection study, its
QLIKE increment is indistinguishable from zero (DM $t = +0.52$); refit daily,
in the same solver, the identical frozen block grows its residual-space
increment to $+0.0108$ and carries it through the raw-scale reconstruction at
$\Delta\text{QLIKE} = -0.00093$, DM $t = +3.56$ --- a result that survived
the study's Romano--Wolf sweep ($p = 0.001$). The mechanism is consistent
with what the products are: predominantly fast pairs built from one-bar
quantities, exactly where stale coefficients cost the most, scored under a
loss whose weight sits on high-volatility bars, exactly where fresh
coefficients matter most. Frozen identity, fast coefficients: which products
matter is a constant of the panel; how much each matters right now is not.

Assembled, the product block is a hundred columns with four documented
properties --- sparse, frozen, heavily shrunk, frequently refit --- and it
converts the tree subsection's specification hint into a concrete object the
paper's estimator can carry.

\subsection{The Three-Block Ridge}
\label{sec:three_block}
% Endpoint: alpha=1 target HAR, alpha=100 linear exog, alpha=1000 product.

The model this section ends on is the two-block ridge of
Section~\ref{sec:dense_weak} with the product block appended as a third
block, each block carrying its own penalty inside one rolling solver:
\begin{equation}
    \hat{y}_t \;=\;
    \underbrace{X^{\mathrm{HAR}}_t \beta_1}_{\alpha_1 = 1}
    \;+\;
    \underbrace{X^{\mathrm{exog}}_t \beta_2}_{\alpha_2 = 100}
    \;+\;
    \underbrace{X^{\mathrm{prod}}_t \beta_3}_{\alpha_3 = 1000},
    \label{eq:three_block}
\end{equation}
where the first block is the target's own HAR ladder, essentially
unpenalized; the second is the linear exogenous block of
Section~\ref{sec:dense_weak} under its heavy uniform shrinkage; and the third
is the frozen product block, shrunk an order of magnitude harder still. The
block structure is imposed by column scaling --- each block's columns scaled
by $\sqrt{\alpha_1/\alpha_b}$ --- so the three-block estimator runs in the
same exact rank-one rolling least-squares path as every other linear model in
this paper, refit at every bar, and nests the two-block model exactly
($\beta_3 = 0$).

The penalty ladder is the section's argument written as three numbers. The
backbone is left alone because Section~\ref{sec:dense_weak} showed that
shrinking a low-dimensional, fully informative block is pure cost. The
exogenous block is shrunk hard because the signal is dense and weak. The
product block is shrunk another decade harder because that is what the
interaction study measured: products live on a different scale from their
parents, and only under heavy block-specific shrinkage is their contribution
positive and stable --- the same one-order-of-magnitude gap between the
linear-block and product-block penalties that the study's own ladder selected
on its panel. Each block's penalty encodes what this paper learned about that
block's information, and no penalty was chosen by grid search on the
evaluation loss.

\pending{the three-block result table on the paper's panel: QLIKE for the
three-block ridge alongside the two-block ridge of Section 5 and the
incumbent, with per-bar DM statistics for the product-block increment ---
the documented analogues (the interaction study's composed models: HAR
backbone per-bar $+$ daily-refit $516$-column exog $+$ $100$ frozen products
at a $10\times$ block-penalty ratio, QLIKE $0.12651$ vs.\ its HAR reference
$0.13275$, DM $t = +13.5$, product increment $t = +3.56$) were computed
under the interaction-study convention ($250$-day window, $27$-column
backbone) and are not quotable as this paper's numbers}
\decide{whether the three-block table, once run, reports the tree arms
alongside it --- the natural closing comparison is the three-block ridge
against the causally tuned LightGBM on \texttt{all\_features} ($0.12604$),
i.e.\ whether naming the tree's structure recovers the tree's level within
the linear class}

What the three-block ridge claims is deliberately modest. It does not claim
to exhaust the trees' nonlinearity --- the functional-ANOVA remainder and the
expressivity ladder both leave a diffuse residue that no small set of named
columns will capture. It claims the constructive converse of this section's
diagnosis: the \emph{recoverable} nonlinearity is pairwise, sparse, frozen,
and clock-anchored, and once it is written down as a block, the estimator
that carries it is still a linear one --- one solver, three penalties, refit
every bar. The next section adds the fourth block.
